\documentclass[11pt]{article}
\usepackage[a4paper,margin=23mm]{geometry}
\usepackage{amsmath,amssymb,amsthm,mathtools,bm}
\usepackage{booktabs,array,microtype,xurl,hyperref}
\hypersetup{colorlinks=true,linkcolor=blue,citecolor=blue,urlcolor=blue}

\newtheorem{theorem}{Theorem}[section]
\newtheorem{proposition}[theorem]{Proposition}
\newtheorem{lemma}[theorem]{Lemma}
\newtheorem{corollary}[theorem]{Corollary}
\newtheorem{definition}[theorem]{Definition}
\newtheorem{remark}[theorem]{Remark}
\newcommand{\dd}{\,\mathrm d}
\newcommand{\RR}{\mathbb R}
\newcommand{\SSS}{\mathbb S}
\newcommand{\tr}{\operatorname{tr}}
\newcommand{\cof}{\operatorname{cof}}
\newcommand{\MM}{\mathcal M}
\newcommand{\GG}{\mathcal G}

\title{AP-E12 Fixed-Degree Complete-Minor Graph Density:\
Exact Strong Quasiconvexity, an Endpoint Dipole No-Go,\
and the Fail-Closed Relaxed-Minimizer Regularity Branch}
\author{Route-F derivation ledger}
\date{20 July 2026}

\begin{document}
\maketitle

\begin{abstract}
AP-E11 left one precise continuum gate: strong smooth density, at fixed
degree, for maps whose first, second, and third derivative minors all belong
to $L^2$.  AP-E12 audits that endpoint instead of assuming it.  Two useful
model-specific theorems are proved.  First, graph-norm convergence of
sphere-valued maps automatically preserves degree, so ``fixed degree'' adds
no second repair once target-valued graph density is known.  Second, the
complete-minor density has an exact strong-quasiconvexity identity and
uniform Legendre--Hadamard constant one.  A scale-sharp calculation also
excludes the direct Mal\'y two-hole/dipole mechanism: finite $L^2$ second-
minor energy requires $\alpha>1/2$, while a non-vanishing forced filling cost
requires $\alpha\leq1/2$, including the logarithmic borderline.

These results do not prove the endpoint density theorem.  The classical
Cartesian approximation theorem loses exponent, and area-approximation
results control all minors only in $L^1$.  The fallback partial-regularity
theory also does not apply: the present density has $(p,q)=(2,6)$ growth in
dimension three, outside the recent relaxed-quasiconvex range $q<3$, and it
has an additional $\mathbb S^3$ constraint and fixed-degree relaxation.  A
finite-grid concentration and perturbed-start card passes $9/9$ and gives a
very strong same-basin proxy after quotienting translations and target
$SO(3)$, but is not a continuum isolation theorem.  Therefore endpoint
density is neither proved nor disproved; partial regularity, classicality,
continuum isolation, the Riemann Hessian, and determinant variation all
remain false in the authorization ledger.
\end{abstract}

\tableofcontents

\section{The exact endpoint statement}

Let $\Omega\subset\RR^3$ be a bounded Lipschitz domain, with constant vacuum
trace $e_0=(1,0,0,0)$, and let $M_k(F)$ denote the vector of all $k\times k$
minors of $F\in\RR^{4\times3}$.  Define the complete-minor graph class
\begin{equation}
 \GG^2_B(\Omega,\SSS^3)=\left\{
 \begin{array}{l|l}
 u\in W^{1,2}(\Omega;\SSS^3)&
 \operatorname{tr}u=e_0,\quad M_2(Du),M_3(Du)\in L^2,\\[-2mm]
 &\deg u=B,\quad M_k(Du)\text{ are distributional minors}
 \end{array}\right\}.
 \label{eq:graph-class}
\end{equation}
The strong graph distance is
\begin{equation}
 d_{\GG}(u,v)=\|u-v\|_{L^2}+\|Du-Dv\|_{L^2}
 +\|M_2(Du)-M_2(Dv)\|_{L^2}
 +\|M_3(Du)-M_3(Dv)\|_{L^2}.
 \label{eq:graph-distance}
\end{equation}
For sphere-valued maps the degree is
\begin{equation}
 B(u)=\frac1{2\pi^2}\int_\Omega
 \det[u,\partial_1u,\partial_2u,\partial_3u]\,\dd x .
 \label{eq:degree}
\end{equation}
It is integrable because $u\in L^\infty$ and $M_3(Du)\in L^2$.

The missing AP-E11 assertion is
\begin{equation}
 \boxed{\overline{C^\infty_{e_0}(\overline\Omega;\SSS^3)
 \cap\{\deg=B\}}^{\ d_{\GG}}=\GG^2_B(\Omega,\SSS^3).}
 \qquad(D_{\rm end})\label{eq:endpoint-density}
\end{equation}
Ordinary $W^{1,2}$ density is available because the relevant obstruction is
$\pi_2(\SSS^3)=0$ \cite{Bethuel1991}; it does not imply
\eqref{eq:endpoint-density}.  In particular, strong $L^2$ convergence of
$Du_j$ gives only $L^1$ control of quadratic products, not strong $L^2$
control of $M_2(Du_j)$, and gives no endpoint control of cubic products.

\begin{proposition}[Degree is automatically closed in the graph norm]
Suppose $u_j,u\in W^{1,2}(\Omega;\SSS^3)$ have the vacuum trace and
$d_{\GG}(u_j,u)\to0$.  Then $B(u_j)\to B(u)$.  If the $u_j$ are smooth, their
integer degrees therefore equal $B(u)$ for all sufficiently large $j$.
\end{proposition}

\begin{proof}
Write $J_3(Du)$ for the signed four-vector dual to $M_3(Du)$.  Since
$\|u_j\|_\infty=\|u\|_\infty=1$,
\begin{align}
 2\pi^2|B(u_j)-B(u)|
 &\leq\int_\Omega |(u_j-u)\cdot J_3(Du_j)|
       +|u\cdot[J_3(Du_j)-J_3(Du)]|\,\dd x\notag\\
 &\leq \|u_j-u\|_2\|M_3(Du_j)\|_2
   +|\Omega|^{1/2}\|M_3(Du_j)-M_3(Du)\|_2.
 \label{eq:degree-continuity}
\end{align}
The right side tends to zero, because the second graph component also bounds
$\|M_3(Du_j)\|_2$.  Integers converging to the integer $B(u)$ are eventually
equal to it.
\end{proof}

Thus the real missing theorem is target-valued endpoint graph density.  A
separate degree-changing patch must not be inserted; if graph convergence is
proved, degree follows from \eqref{eq:degree-continuity}.

\section{Exact algebra of the AP-E11 continuum density}

For $R,K>0$ put
\begin{equation}
 W(F)=\frac12\left(|F|^2+R|M_2(F)|^2+K|M_3(F)|^2\right).
 \label{eq:density}
\end{equation}
If $s_1,s_2,s_3$ are the singular values of $F$, Cauchy--Binet gives
\begin{equation}
 2W(F)=\sum_i s_i^2+R\sum_{i<j}s_i^2s_j^2
       +K s_1^2s_2^2s_3^2.
 \label{eq:singular}
\end{equation}
This is convex in the complete-minor vector $(F,M_2(F),M_3(F))$, but not in
$F$ itself.

\begin{theorem}[Uniform rank-one symbol]
For every $F\in\RR^{4\times3}$ and every rank-one
$H=a\otimes\xi$,
\begin{equation}
 D^2W(F)[H,H]=|H|^2+R|DM_2(F)[H]|^2+K|DM_3(F)[H]|^2
 \geq |a|^2|\xi|^2.
 \label{eq:rank-one}
\end{equation}
Hence the Legendre--Hadamard constant is one, independently of $F$.
\end{theorem}

\begin{proof}
Every minor of $F+t(a\otimes\xi)$ is affine in $t$: a determinant cannot
contain two columns from the same rank-one update.  Taking two derivatives of
the squared Euclidean norms in \eqref{eq:density} gives
\eqref{eq:rank-one}.  Finally $|a\otimes\xi|^2=|a|^2|\xi|^2$.
\end{proof}

\begin{theorem}[Exact strong quasiconvexity identity]
Let $U$ be a periodic cell, $F$ constant, and
$\varphi\in W^{1,\infty}_{\rm per}(U;\RR^4)$ have zero mean.  Then
\begin{align}
 &\int_U\{W(F+D\varphi)-W(F)\}\,\dd x\notag\\
 &\quad=\frac12\int_U\bigl\{|D\varphi|^2
 +R|M_2(F+D\varphi)-M_2(F)|^2
 +K|M_3(F+D\varphi)-M_3(F)|^2\bigr\}\,\dd x.
 \label{eq:strong-qc}
\end{align}
The same formula holds for zero-boundary variations on a bounded domain.
\end{theorem}

\begin{proof}
Every first, second, and third minor is a null Lagrangian.  Periodicity gives
\begin{equation}
 \frac1{|U|}\int_U D\varphi=0,\qquad
 \frac1{|U|}\int_U M_k(F+D\varphi)=M_k(F),\quad k=2,3.
 \label{eq:null-means}
\end{equation}
Expand each squared norm in \eqref{eq:density}.  All cross terms integrate to
zero by \eqref{eq:null-means}; the remaining squares are precisely
\eqref{eq:strong-qc}.
\end{proof}

Equation \eqref{eq:strong-qc} is stronger than a generic quasiconvexity
estimate and is the main positive analytic structure discovered in AP-E12.
It controls complete-minor excess whenever one can construct a competitor
with the correct trace.  It does not itself construct the trace-preserving,
target-valued endpoint recovery sequence needed in
\eqref{eq:endpoint-density}.

\section{What the Cartesian approximation theorems do and do not give}

Mal\'y's approximation theorem for Cartesian maps gives strong convergence
of all minors in exponents $q_k<p_k$; it explicitly does not claim a bounded
approximating sequence at the endpoint \cite{Maly1991}.  Mucci's slicing
characterization treats strong area approximation, equivalently
$W^{1,1}$ plus $L^1$ convergence of all minors \cite{Mucci2001}.  The
Cartesian-current framework supplies the correct closure language for maps
into spheres \cite{GiaquintaModicaSoucek1989}.  None of these statements is
the $L^2$ endpoint \eqref{eq:endpoint-density}.

There is a useful conditional implication.  If the selected map belongs to
the appropriate target-valued Cartesian closure, its complete-minor vector
has local $L^{2+\delta}$ control for some $\delta>0$, and the boundary/
$\mathbb S^3$ projection step can be performed without losing that control,
then Mal\'y's exponent-loss theorem may be applied with output exponent two.
This reduces one route to the higher-integrability estimate
\begin{equation}
 (Du,M_2(Du),M_3(Du))\in L^{2+\delta}_{\rm loc}.
 \qquad(H_\delta)\label{eq:higher-int}
\end{equation}
The target projection hypothesis is essential: ambient smooth approximation
alone does not keep values on $\mathbb S^3$.

\section{A scale-sharp dipole audit}

The natural counterexample attempt is a shrinking two-dimensional
Cartesian dipole extended along a third coordinate.  In a target chart write
\begin{equation}
 u(r,\theta,z)=a(r)\gamma(\theta),
 \label{eq:dipole}
\end{equation}
where $\gamma$ carries the nontrivial two-hole filling datum.  Up to harmless
angular constants,
\begin{align}
 \int |Du|^2&\sim\int_0 a'(r)^2r+\frac{a(r)^2}{r}\,\dd r,
 \label{eq:dipole-d1}\\
 \int |M_2(Du)|^2&\sim\int_0\frac{a(r)^2a'(r)^2}{r}\,\dd r.
 \label{eq:dipole-m2}
\end{align}

\begin{lemma}[No power-law overlap]
For $a(r)=r^\alpha$, finite Dirichlet energy requires $\alpha>0$, while
finite $L^2$ second-minor energy requires $\alpha>1/2$.  Any smoothing that
fills a fixed nonzero signed target area on a disk of radius $r$ obeys
\begin{equation}
 \int_{D_r}|\det Dv|^2\,\dd x
 \geq\frac1{\pi r^2}\left|\int_{D_r}\det Dv\,\dd x\right|^2
 \gtrsim\frac{a(r)^4}{r^2}=r^{4\alpha-2}.
 \label{eq:fill-lower}
\end{equation}
Thus a non-vanishing or divergent endpoint filling cost requires
$\alpha\leq1/2$, which has no overlap with \eqref{eq:dipole-m2}.
\end{lemma}

\begin{proof}
Substitution into \eqref{eq:dipole-d1} gives an integrand proportional to
$r^{2\alpha-1}$; substitution into \eqref{eq:dipole-m2} gives
$r^{4\alpha-3}$.  The stated integrability thresholds follow.  Inequality
\eqref{eq:fill-lower} is Cauchy--Schwarz on $D_r$, and the signed filling area
scales as $a(r)^2$.
\end{proof}

The logarithmic borderline also fails to open a gap.  For
\begin{equation}
 a(r)=r^{1/2}[\log(e/r)]^{-\beta},
\end{equation}
the second-minor integral is finite exactly when $\beta>1/4$, whereas the
lower bound in \eqref{eq:fill-lower} is
$[\log(e/r)]^{-4\beta}\to0$ for every $\beta>0$.  Consequently:
\begin{equation}
 \boxed{\text{the direct self-similar Mal\'y dipole mechanism cannot refute
 the AP-E12 endpoint.}}
 \label{eq:dipole-no-go}
\end{equation}
This is a mechanism-specific no-go, not a proof that no more anisotropic or
iterated concentration construction exists.

\section{Finite energy does not imply continuity}

The regularity fallback cannot begin by assuming that every graph-space map
is continuous.  Around a spatial axis set
\begin{equation}
 u(r,z)=\bigl(\cos\ell(r),\sin\ell(r),0,0\bigr),\qquad
 \ell(r)=\log\log(e/r).
 \label{eq:rank-one-discontinuous}
\end{equation}
Then $Du$ has rank one, so $M_2(Du)=M_3(Du)=0$, while
\begin{equation}
 \int_0^{r_0}|\partial_ru|^2r\,\dd r
 =\int_0^{r_0}\frac{\dd r}{r\log^2(e/r)}<\infty.
 \label{eq:rank-one-energy}
\end{equation}
The phase has no limit as $r\downarrow0$, so there is no continuous
representative on the axis.  Nevertheless, making $\ell$ constant below
$\epsilon$ produces smooth rank-one truncations with zero higher minors and
$W^{1,2}$ tail energy tending to zero.  This example is therefore not a
density counterexample.  It proves only that results assuming continuous
Cartesian maps cannot cover the full energy class without another argument.

\section{The relaxed-minimizer fallback and its exact obstruction}

Add the lower-order mass potential and define
\begin{equation}
 E_B(u)=\int_\Omega\{W(Du)+m^2(1-u^0)\}\,\dd x,\qquad u\in\GG_B^2,
 \label{eq:energy}
\end{equation}
and let $\overline E_B$ be its lower-semicontinuous relaxation from smooth
fixed-degree maps.  AP-E11 proves convergence to minimizers of
$\overline E_B$, not equality $\overline E_B=E_B$ on all of $\GG_B^2$.

For a classical map, put $F=Du$ and $G=F^TF$.  Cauchy--Binet gives
$|M_3(F)|^2=\det G$ and
$|M_2(F)|^2=((\tr G)^2-\tr G^2)/2$.  Hence the exact first Piola stress is
\begin{equation}
 P(F)=DW(F)=F+R\{(\tr G)F-FG\}+KF\cof G.
 \label{eq:stress}
\end{equation}
The formal tangent Euler--Lagrange system is
\begin{equation}
 (I-u\otimes u)\{-\operatorname{div}P(Du)-m^2e_0\}=0.
 \label{eq:projected-el}
\end{equation}

\begin{proposition}[Dependency chain for the fallback]
A global minimizer of $\overline E_B$ may be promoted to a classical isolated
solution of \eqref{eq:projected-el} only after the following independent
steps are supplied:
\begin{enumerate}
 \item local target-valued recovery with fixed trace, identifying relaxed
       compactly supported variations with variations of $E_B$;
 \item a weak Euler--Lagrange theorem for the resulting graph-space local
       minimizer;
 \item an $\epsilon$-regularity or partial-regularity theorem compatible with
       the $\mathbb S^3$ constraint and complete-minor energy;
 \item removal of the possible singular set and a strict local growth
       estimate modulo translations and the global target $SO(3)$ orbit.
\end{enumerate}
Without item 1, minimality of the relaxation does not logically imply
\eqref{eq:projected-el} for the naive integral.
\end{proposition}

The recent theorem of Gmeineder and Kristensen proves smooth partial
regularity of relaxed minimizers for strongly quasiconvex $(p,q)$ integrands
in the range
\begin{equation}
 1\leq p\leq q<\min\left\{\frac{np}{n-1},p+1\right\}
 \cite{GmeinederKristensen2024}.
 \label{eq:gk-range}
\end{equation}
Here the pointwise bounds are $c|F|^2\lesssim W(F)\lesssim C(1+|F|^6)$, so
$(n,p,q)=(3,2,6)$ and \eqref{eq:gk-range} would require $q<3$.
Sch\"affner's higher-integrability theorem for convex nonstandard-growth
integrands requires $q/p<1+2/(n-1)=2$, whereas the present ratio is three
and $W$ is polyconvex rather than convex in $F$ \cite{Schaffner2021}.  Neither
theorem includes the sphere constraint or fixed-degree graph relaxation.
The exact identity \eqref{eq:strong-qc} is promising additional structure,
but applying it requires the missing trace-preserving graph competitor.

Consequently the requested fallback has been executed as far as presently
justified and stops before its first promotion:
\begin{equation}
 \boxed{C_{\rm local\ recovery}=C_{\rm weak\ EL}
 =C_{\rm partial\ regularity}=C_{\rm classicality}
 =C_{\rm isolation}=\mathrm{false}.}
 \label{eq:fallback-false}
\end{equation}
This is an absence of proof, not a counterexample to regularity of the
particular numerical $B=1$ minimizer.

\section{Numerical concentration and isolation proxy}

The production script re-relaxes the same AP-E11 action at three spacings and
volumes.  No center or core value is pinned.  The maximum fraction of total
energy in one tetrahedron decreases under refinement, as does the $99.9$th
percentile of the local energy density:
\begin{center}
\begin{tabular}{@{}rrrrrr@{}}
\toprule
$N$&$L$&$a$&$E$&max tet fraction&density $q_{0.999}$\\
\midrule
17&5.5&0.343750&77.5780&$3.966\times10^{-3}$&31.187\\
21&6.0&0.300000&79.2457&$1.849\times10^{-3}$&24.900\\
25&6.5&0.270833&79.9500&$1.173\times10^{-3}$&20.971\\
\bottomrule
\end{tabular}
\end{center}
All three have $B=(1,1,1)$, minimum pair dot at least $0.2724$, and projected
gradient density at most $1.54\times10^{-7}$.  In the smallest audited balls,
whose radii are approximately two lattice spacings, the energy fractions
decrease $0.420\to0.264\to0.189$.  These observations exclude a visible
one-cell concentration in this range; they are not a compensated-compactness
or $\epsilon$-regularity theorem.

Three independent smooth tangent perturbations of amplitude $0.24$ are then
re-relaxed on the $N=17$ action.  Field distance is evaluated after quotienting
integer spatial translations and the unbroken target $SO(3)$.  If the pion
parts of reference and candidate are matrices $X,Y$, the proper Procrustes
rotation is
\begin{equation}
 Y^TX=U\Sigma V^T,\qquad
 Q_*=U\operatorname{diag}(1,1,\det(UV^T))V^T\in SO(3).
 \label{eq:procrustes}
\end{equation}
The maximum results over the three starts are
\begin{equation}
 \frac{\max E-\min E}{\overline E}=3.38\times10^{-11},\qquad
 d_{\rm CDF}=5.70\times10^{-7},\qquad
 d_{\rm field/(trans\times SO(3))}=1.31\times10^{-5}.
 \label{eq:finite-isolation}
\end{equation}
All endpoints remain $B=(1,1,1)$ and stationary below
$4.31\times10^{-7}$.  The predeclared finite-dimensional basin proxy passes.
It is deliberately not promoted to continuum isolation because it samples
only three perturbations on one finite grid and does not give a quantitative
neighbourhood theorem.

\section{Decision and next theorem}

The reproducible card
\begin{center}
\texttt{route\_f/code/verify\_ap\_e12\_graph\_density\_regular\_minimizer.py}
\end{center}
passes $9/9$ algebraic, scaling, quadrature, concentration, and finite-basin
checks.  The logical authorization is nevertheless
\begin{equation}
 \boxed{\begin{gathered}
 C_{\rm degree\ closed\ in\ graph}=C_{\rm LH}
 =C_{\rm strong\ QC}=C_{\rm dipole\ no\mbox{-}go}=\mathrm{true},\\
 C_{D_{\rm end}}=\mathrm{false\ (not\ disproved)},\\
 C_{\rm partial\ regularity}=C_{\rm classicality}=C_{\rm isolation}
 =C_{\rm Hessian}=C_{\det}=\mathrm{false}.
 \end{gathered}}
 \label{eq:decision}
\end{equation}

The most targeted next lemma is an annular, target-valued endpoint replacement:
given a low-energy shell, choose a good sphere and replace its interior while
preserving the trace and obtaining equiintegrability of all three $L^2$ minor
families.  Combined with \eqref{eq:strong-qc}, such a lemma would yield either
a graph-density proof or the Caccioppoli inequality needed for the fallback.
The second route is to prove \eqref{eq:higher-int} directly for the selected
relaxed minimizer and then establish the sphere-valued projection/retraction
step.  Only after one route proves classicality and continuum isolation is it
legal to assemble the same-action Riemann Hessian.  The general danger that a
relaxed and a smooth variational problem need not have the same infimum is
precisely the Lavrentiev boundary emphasized in recent reviews
\cite{CerfMariconda2024}; it must not be removed by numerical agreement alone.

\bibliographystyle{unsrt}
\bibliography{route_f/tex/ap_e12_graph_density_regular_minimizer}

\end{document}
